% =====================================================================
%  Avant-propos
% =====================================================================
\chapter*{Avant-propos}
\addcontentsline{toc}{chapter}{Avant-propos}
\markboth{Avant-propos}{Avant-propos}

\section*{Pourquoi ce module ?}

Django est l'un des frameworks web les plus utilisés au monde. Il permet de construire
rapidement des applications solides, sécurisées et maintenables avec Python. Pourtant,
les débutants se perdent souvent dans la documentation officielle, très complète mais
peu progressive. Ce module a été écrit pour combler cet écart.

Nous avons choisi un fil rouge unique : une petite application de gestion de
bibliothèque appelée \gesibiblio. Chaque chapitre ajoute une brique à cette
application. À la fin, vous ne disposez pas seulement de notes de cours : vous avez
un projet fonctionnel, que vous pouvez montrer, expliquer et faire évoluer.

\section*{À qui s'adresse ce module ?}

Ce module est destiné aux personnes qui débutent en développement web et qui
souhaitent apprendre Django par la pratique. Il ne suppose aucune connaissance
préalable de Django. En revanche, une familiarité minimale avec Python, le terminal
et le HTML est un vrai atout. Les rappels nécessaires sont fournis au fil des
chapitres et dans les annexes.

\begin{objectifs}
À la fin de ce module, vous serez capable de :
\begin{itemize}
  \item installer Python, créer un environnement virtuel et installer Django ;
  \item expliquer le rôle de chaque fichier d'un projet Django ;
  \item créer une application, des vues, des URL et des gabarits ;
  \item concevoir des modèles de données et maîtriser les migrations ;
  \item interroger la base avec l'ORM de Django ;
  \item gérer les requêtes HTTP, les formulaires et les fichiers envoyés ;
  \item utiliser l'administration Django et personnaliser son affichage ;
  \item mettre en place l'authentification, les groupes et les permissions ;
  \item appliquer les bonnes pratiques de structuration et de sécurité.
\end{itemize}
\end{objectifs}

\section*{La méthode : apprendre en construisant}

Chaque chapitre suit la même progression :
\begin{enumerate}
  \item une situation concrète, racontée simplement ;
  \item le concept Django correspondant, expliqué sans jargon inutile ;
  \item la mise en pratique dans le projet \gesibiblio, avec le code complet ;
  \item un encadré \og À retenir \fg{} qui fixe les idées clés ;
  \item des erreurs fréquentes, pour ne pas perdre de temps ;
  \item un ou deux exercices, avec une correction détaillée.
\end{enumerate}

Le code présenté est du code réel, testé avec Django 5.2 LTS. Les noms de fichiers,
les variables et les URL sont cohérents d'un chapitre à l'autre. Vous pouvez donc
taper le code au fur et à mesure, sans avoir à réinventer l'architecture.

\begin{bonne}
Tapez le code vous-même. La lecture passive donne l'illusion de comprendre. La
frappe, les erreurs de syntaxe et les messages du terminal font partie de
l'apprentissage. Prenez le temps de lire chaque message d'erreur : Django est
généralement très explicite.
\end{bonne}

\section*{Le plan de la formation}

\begin{table}[htbp]
\centering
\small
\begin{tabularx}{\textwidth}{C{1.3cm} L{8.4cm} C{1.7cm}}
\toprule
\textbf{\#} & \textbf{Contenu} & \textbf{Durée} \\
\midrule
1 & Introduction : le Web, Django, le modèle MVT et l'environnement & 3 h \\
2 & Structure d'un projet Django : rôles de chaque fichier, réglages, serveur & 4 h \\
3 & Création d'une application : vues, URL et première page & 3 h \\
4 & Chargement d'un gabarit : héritage, blocs, inclusions & 4 h \\
5 & Fichiers statiques : CSS, images, JavaScript et fichiers média & 4 h \\
6 & Balises et filtres : affichage dynamique, filtres personnalisés & 4 h \\
7 & Modèles de données : ORM, relations, migrations et requêtes & 6 h \\
8 & Gestion des requêtes : HTTP, vues, recherche, pagination & 4 h \\
9 & Administration Django : inscription des modèles et personnalisation & 3 h \\
10 & Formulaires Django : validation, création, modification, suppression & 5 h \\
11 & Authentification : comptes, connexion, profil et sessions & 3 h \\
12 & Permissions : groupes, droits, sécurité et bonnes pratiques & 3 h \\
13 & Conclusion : bilan, tests, déploiement et suite du parcours & 2 h \\
\midrule
 & \textbf{Total} & \textbf{48 h} \\
\bottomrule
\end{tabularx}
\caption{Plan détaillé du module. La durée inclut les démonstrations et les exercices.}
\end{table}

Le volume indiqué correspond à une formation guidée. En autonomie, comptez environ
le double, surtout si vous refaites chaque exercice sans regarder la correction.

\section*{Prérequis et matériel}

\begin{itemize}
  \item Un ordinateur sous Windows, macOS ou Linux, avec au moins 4 Go de RAM.
  \item Python 3.10 ou plus récent. Nous utilisons Python 3.12 dans les exemples.
  \item Un éditeur de code. Visual Studio Code est recommandé, mais PyCharm ou
        Sublime Text conviennent très bien.
  \item Une connexion Internet pour installer les paquets et consulter la
        documentation.
  \item Aucun serveur web, aucune base de données à installer : Django fournit un
        serveur de développement et utilise SQLite par défaut.
\end{itemize}

\section*{Conventions utilisées}

\begin{description}[leftmargin=3.2cm,style=nextline]
  \item[Code Python] Les blocs avec un fond gris clair et des numéros de ligne
  contiennent du code à recopier. Le nom du fichier est indiqué dans le bandeau.
  \item[Terminal] Les blocs sombres contiennent des commandes à exécuter dans le
  terminal. Le signe \code{\$} n'est pas à recopier sous Windows.
  \item[Encadrés] Les encadrés \og À retenir \fg, \og Bonne pratique \fg,
  \og Attention \fg, \og Erreur fréquente \fg{} et \og Note \fg{} mettent en avant
  les points importants.
  \item[Exercices] Les exercices sont numérotés par chapitre. Une correction
  détaillée suit immédiatement, dans un encadré gris.
  \item[Chemins] Les chemins sont donnés avec des barres obliques, comme dans le
  code Python. Sous Windows, Django accepte aussi les barres obliques.
\end{description}

\section*{Versions et compatibilité}

Le module a été rédigé pour \textbf{Django 5.2 LTS} et \textbf{Python 3.12}. Django
5.2 est une version de support à long terme, ce qui est un choix raisonnable pour
un projet de formation. Les concepts présentés restent valables pour Django 4.2 et
pour les versions 5.x suivantes. Les différences éventuelles sont signalées dans
le texte.

\begin{attention}
Si vous utilisez une version plus ancienne de Django, certains détails peuvent
changer, notamment la syntaxe de \code{CheckConstraint} ou le comportement de
\code{LogoutView}. Vérifiez la version avec \code{python -m django --version} et
privilégiez Django 5.2 LTS pour suivre le module sans adaptation.
\end{attention}

\section*{Remerciements}

Ce module a été conçu pour la formation \og Initiation au développement web avec
Django \fg, propulsée par \textbf{GESITECHNOLOGIES}. Merci aux apprenants et aux
formateurs qui ont testé les exercices et proposé des améliorations. Leur regard
concret a rendu ce document plus utile.

\vspace{1cm}
\begin{flushright}
\textit{LWANZO PATERNE}\\
\textit{Formateur en développement web}
\end{flushright}
